\section{Panorama de la serie e introducción del curso}

Esta entrega es la primera de la serie de paga Modern Agents. Desde una perspectiva macro, hilvana los temas centrales que la serie abordará en el futuro y, tomando \textbf{Reasoning \& Planning} como punto de partida, presenta un survey sistemático de las capacidades y las carencias de los modelos de lenguaje de gran escala (LLM) en el ámbito de la planeación.

\subsection{Temas que cubrirá la serie}

La serie Modern Agents en su conjunto planea cubrir los siguientes grandes bloques:

\begin{enumerate}
    \item \textbf{Reasoning and Planning}: la capacidad de razonamiento y planeación de los modelos de lenguaje, que es la base de todo el contenido posterior.
    \item \textbf{Workflows clásicos}: incluyen Self-Critique (el ciclo de generación, evaluación e iteración) y Multi-Agent System (la colaboración entre múltiples agentes, cuyo núcleo es la Separation of Concern).
    \item \textbf{Context Engineering}: el empaquetado integral de las técnicas relacionadas con los modelos de lenguaje, con énfasis en el diseño de Memory.
    \item \textbf{Self-Evolving}: técnicas de automejora representadas por AlphaEvolve, GPA y similares.
    \item \textbf{Práctica con frameworks}: el uso de frameworks como LangChain / LangGraph.
    \item \textbf{Métodos basados en grafos}: soluciones Neuro-Symbolic como GraphRAG.
    \item \textbf{Diseño y encapsulado de herramientas}: diseñar Tools / Servers adecuados para los modelos de lenguaje, de modo que compensen sus carencias.
    \item \textbf{Test-Time Computing}: técnicas de cómputo en tiempo de inferencia, como el pensamiento en paralelo (DeepConf) y el Majority Voting.
\end{enumerate}

\subsection{Resumen de la sección}

La idea central de esta serie es: dominar, mediante la práctica continua, cada vez más know-how eficaz, y maximizar el desempeño de los modelos de lenguaje a través del diseño y el desarrollo de Agents.

